% Jerarquía con retroalimentación: cada nivel sube del anterior (`derivacion`) y el nivel
% superior reelabora los inferiores con un haz de arcos punteados (`retroalim`).
\documentclass[tikz,border=12pt]{standalone}
\usepackage{diagrama}
\begin{document}
\begin{tikzpicture}[every node/.append style={font=\px{26}}]
  \node[titulo=40] at (3,6.8) {Título de la jerarquía};
  \foreach \i/\y/\t in {1/0.6/Nivel 1, 2/1.9/Nivel 2, 3/3.2/Nivel 3, 4/4.5/Nivel 4}
    \node[caja, minimum width=2.6cm, minimum height=0.9cm] (n\i) at (3,\y) {\t};
  \foreach \i/\j in {1/2, 2/3, 3/4}
    \draw[derivacion] (n\i.north) -- (n\j.south);          % e\i
  % el nivel superior reelabora cada uno de los inferiores
  \draw[retroalim] (n4.west) to[out=180, in=180, looseness=1.6] (n1.west);   % e4
  \draw[retroalim] (n4.west) to[out=180, in=180, looseness=1.2] (n2.west);   % e5
  \draw[retroalim] (n4.west) to[out=180, in=180, looseness=0.9] (n3.west);   % e6
  \node[font=\itshape\px{20}] at (0.1,3.2) {Explicitación\\global};
  \draw[eje virtual] (5.6,0.4) -- node[rotulo=20, rotate=90] {\textit{Incrustación funcional}} (5.6,4.7);  % x1
\end{tikzpicture}
\end{document}
